\documentclass[11pt,a4paper]{article}
\usepackage[spanish,es-noquoting]{babel}
\usepackage[top=2.2cm, bottom=2.2cm, left=2.5cm, right=2.5cm]{geometry}
\usepackage{fontspec}
\usepackage{amsmath}
\usepackage[version=4]{mhchem}
\usepackage{enumitem}
\usepackage{microtype}

\setmainfont{Calibri}
\linespread{1.15}

\begin{document}
\pagestyle{empty}

\noindent\textbf{1. Resolución del ejercicio de síntesis de amoníaco (\ce{NH3}):}

\vspace{0.2cm}
\noindent\textbf{Cálculo de masas molares ($M$):}
\begin{itemize}[leftmargin=1.5em, itemsep=2pt]
    \item $M(\ce{N2}) = 2 \cdot 14 = 28\text{ g/mol}$
    \item $M(\ce{H2}) = 2 \cdot 1 = 2\text{ g/mol}$
    \item $M(\ce{NH3}) = 14 + 3(1) = 17\text{ g/mol}$
\end{itemize}

\vspace{0.15cm}
\noindent\textbf{a) Ajuste de la reacción:}
\[
\ce{N2 (g) + 3 H2 (g) -> 2 NH3 (g)}
\]

\vspace{0.1cm}
\noindent\textbf{b) Reactivo limitante y reactivo en exceso:}
\begin{itemize}[leftmargin=1.5em, itemsep=2pt]
    \item Cantidades iniciales en moles:
    \[
    n(\ce{N2}) = \frac{140\text{ g}}{28\text{ g/mol}} = 5{,}00\text{ mol de \ce{N2}} \qquad n(\ce{H2}) = \frac{36\text{ g}}{2\text{ g/mol}} = 18{,}0\text{ mol de \ce{H2}}
    \]
    \item Según la estequiometría ($1\text{ mol \ce{N2}} : 3\text{ mol \ce{H2}}$), para consumir $5{,}00\text{ mol de \ce{N2}}$:
    \[
    n(\ce{H2})_{\text{necesario}} = 5{,}00\text{ mol \ce{N2}} \times \frac{3\text{ mol \ce{H2}}}{1\text{ mol \ce{N2}}} = 15{,}0\text{ mol de \ce{H2}}
    \]
    \item Como disponemos de $18{,}0\text{ mol de \ce{H2}}$ ($18{,}0 > 15{,}0$), el \textbf{reactivo limitante} es el nitrógeno (\ce{N2}) y el \textbf{reactivo en exceso} es el hidrógeno (\ce{H2}).
\end{itemize}

\vspace{0.1cm}
\noindent\textbf{c) Masa teórica máxima de \ce{NH3} ($100\%$ de rendimiento):}
\[
n_{\text{teórico}}(\ce{NH3}) = 5{,}00\text{ mol \ce{N2}} \times \frac{2\text{ mol \ce{NH3}}}{1\text{ mol \ce{N2}}} = 10{,}0\text{ mol de \ce{NH3}}
\]
\[
m_{\text{teórica}}(\ce{NH3}) = 10{,}0\text{ mol} \times 17\text{ g/mol} = \mathbf{170\text{ g de \ce{NH3}}}
\]

\vspace{0.1cm}
\noindent\textbf{d) Masa real obtenida ($80\%$ de rendimiento):}
\[
m_{\text{real}}(\ce{NH3}) = m_{\text{teórica}} \times 0{,}80 = 170\text{ g} \times 0{,}80 = \mathbf{136\text{ g de \ce{NH3}}}
\]

\vspace{0.1cm}
\noindent\textbf{e) Reactivo en exceso sobrante:}
\begin{itemize}[leftmargin=1.5em, itemsep=2pt]
    \item Moles sobrantes: $n_{\text{sobrante}} = 18{,}0\text{ mol} - 15{,}0\text{ mol} = \mathbf{3{,}0\text{ mol de \ce{H2}}}$.
    \item Masa sobrante: $m_{\text{sobrante}} = 3{,}0\text{ mol} \times 2\text{ g/mol} = \mathbf{6{,}0\text{ g de \ce{H2}}}$.
\end{itemize}

\newpage

\noindent\textbf{2. Resolución del ejercicio de descomposición térmica de \ce{CaCO3}:}

\vspace{0.2cm}
\noindent\textbf{Cálculo de masas molares ($M$):}
\begin{itemize}[leftmargin=1.5em, itemsep=2pt]
    \item $M(\ce{CaCO3}) = 40 + 12 + 3(16) = 100\text{ g/mol}$
    \item $M(\ce{CaO}) = 40 + 16 = 56\text{ g/mol}$
    \item $M(\ce{CO2}) = 12 + 2(16) = 44\text{ g/mol}$
\end{itemize}

\vspace{0.15cm}
\noindent\textbf{a) Ajuste de la reacción química:}
\[
\ce{CaCO3 (s) -> CaO (s) + CO2 (g)} \quad (\text{Relación estequiométrica } 1:1:1)
\]

\vspace{0.1cm}
\noindent\textbf{b) Masa teórica máxima de óxido de calcio (\ce{CaO}):}
\begin{itemize}[leftmargin=1.5em, itemsep=2pt]
    \item Moles iniciales de reactivo:
    \[
    n(\ce{CaCO3}) = \frac{250\text{ g}}{100\text{ g/mol}} = 2{,}50\text{ mol de \ce{CaCO3}}
    \]
    \item Moles teóricos y masa teórica de \ce{CaO} producida:
    \[
    n_{\text{teórico}}(\ce{CaO}) = 2{,}50\text{ mol \ce{CaCO3}} \times \frac{1\text{ mol \ce{CaO}}}{1\text{ mol \ce{CaCO3}}} = 2{,}50\text{ mol de \ce{CaO}}
    \]
    \[
    m_{\text{teórica}}(\ce{CaO}) = 2{,}50\text{ mol} \times 56\text{ g/mol} = \mathbf{140\text{ g de \ce{CaO}}}
    \]
\end{itemize}

\vspace{0.1cm}
\noindent\textbf{c) Rendimiento porcentual del proceso:}
\[
\text{Rendimiento} = \frac{m_{\text{real}}}{m_{\text{teórica}}} \times 100 = \frac{112\text{ g}}{140\text{ g}} \times 100 = \mathbf{80\%}
\]

\vspace{0.1cm}
\noindent\textbf{d) Masa real de dióxido de carbono (\ce{CO2}) desprendida:}
\begin{itemize}[leftmargin=1.5em, itemsep=2pt]
    \item Cantidad teórica esperada de \ce{CO2}:
    \[
    n_{\text{teórico}}(\ce{CO2}) = 2{,}50\text{ mol \ce{CaCO3}} \times \frac{1\text{ mol \ce{CO2}}}{1\text{ mol \ce{CaCO3}}} = 2{,}50\text{ mol de \ce{CO2}}
    \]
    \[
    m_{\text{teórica}}(\ce{CO2}) = 2{,}50\text{ mol} \times 44\text{ g/mol} = 110\text{ g de \ce{CO2}}
    \]
    \item Masa real obtenida aplicando el rendimiento del $80\%$:
    \[
    m_{\text{real}}(\ce{CO2}) = m_{\text{teórica}} \times 0{,}80 = 110\text{ g} \times 0{,}80 = \mathbf{88\text{ g de \ce{CO2}}}
    \]
\end{itemize}

\end{document}
